\section{Specialization over a Higgs line}\label{sec:specialization}

We now prove the perverse-filtration bound for the operator $f'$ in
\eqref{trans:fprime}.  The translation is defined over a logarithmic Hitchin
base away from zero on a Higgs line.  Nearby cycles extend its action to the
zero fiber as a sheaf endomorphism.  A final integration over the moving-point
curve supplies the shift by $-2$.  The essential comparison with ordinary
cohomology uses the restriction isomorphisms of
Corollary~\ref{log:restriction}; it does not require the parameter morphism
to be proper.

\subsection{The proper morphism and its relative translation}

Keep the vector $v$ and the translation from Section~\ref{sec:translation}.
Let $p_C:C\times S\to S$ be the projection and let $\Delta\subset C\times S$
be the graph of $x:S\to C$.  The vector bundle of trace-free logarithmic
Hitchin coefficients is
\[
 \mathcal E=
 \bigoplus_{i=2}^n
 p_{C,*}\bigl(K_{C\times S/S}(\Delta)^{\otimes i}\bigr).
\]
We use the same notation for its total space.  These direct images are locally
free and commute with base change: each twisting line on a fiber has degree
$i(2g-1)>2g-2$, so its first cohomology vanishes.  The inclusion of ordinary
coefficient spaces defines a closed vector-subbundle embedding
\[
 j:S\times A\hookrightarrow\mathcal E.
\]
Indeed both the ordinary spaces and the logarithmic spaces commute with base
change, and their quotient has constant rank.

Let $D=\mathbb A^1_z$, let $D^*=D\setminus\{0\}$, and define
\[
 \mathcal Y_v=\mathcal Y\times_{\Lambda}D,
 \qquad z\longmapsto(0,zv).
\]
Thus its fiber at $z$ is $Y_{0,zv}$, including the moving-point coordinate in
$S$.  Characteristic coefficients and the parameter define
\begin{equation}\label{spec:hitchin-line}
 g_v:\mathcal Y_v\longrightarrow\mathcal E\times D.
\end{equation}
Write $g_0:Y_{0,0}\to\mathcal E$ for its zero fiber.

\begin{lemma}\label{spec:geometry}
The morphism $g_v$ is proper, and $\mathcal Y_v\to D$ is smooth.  The
translation defines an automorphism $T_v$ of $\mathcal Y_v|_{D^*}$ over
$\mathcal E\times D^*$.  Scaling gives an isomorphism
\[
 \mathcal Y_v|_{D^*}\simeq Y_{0,v}\times D^*
\]
over $D^*$, and restriction from $H^*(\mathcal Y_v)$ to either $Y_{0,0}$ or
$Y_{0,v}$ is an isomorphism.
\end{lemma}

\begin{proof}
Smoothness and the restriction assertions follow from
Proposition~\ref{log:families} and Corollary~\ref{log:restriction}.  For properness,
start with the relative logarithmic Higgs Hitchin morphism over $S$.
Imposing determinant $\mathcal O_C(dx(s))$ and trace zero gives a closed
subspace.  Adding a residue-invariant full flag is projective over that
subspace.  Finally the requirement that the ordered diagonal residues equal
$zv_i$ is closed after adjoining $z$.  The properness of the logarithmic
Hitchin morphism from Section~\ref{sec:logarithmic} therefore implies
properness of \eqref{spec:hitchin-line}.  Coprimality ensures that the
semistable limits in this argument remain stable.

For $z\ne0$ the residual eigenvalues $zv_i$ are distinct, so the elementary
modifications defining $T_v$ and their inverses are algebraic in $z$.
They remain in the stable moduli by the trace-residue argument of
Section~\ref{sec:translation}.  A modified Higgs field agrees with the
original meromorphic Higgs field away from the pole.  Their characteristic
coefficients are consequently equal as sections of
$K_C(x(s))^{\otimes i}$, which proves that $T_v$ is over $\mathcal E$.
Multiplying the Higgs field by $z$ identifies the fiber at $1$ with that at
$z$; division by $z$ is its inverse.  This gives the claimed product over
$D^*$.  The product need not fix the Hitchin coefficient in $\mathcal E$;
only the translation must do so.
\end{proof}

\subsection{Specialization and ordinary cohomology}

Here is the comparison needed for the noncompact spaces in
Lemma~\ref{spec:geometry}.  Throughout, nearby cycles are unshifted;
we use the nearby-cycle and proper-base-change formalism recalled in
\cite[Sections~5.5 and~5.8]{dCM2009}.

\begin{lemma}\label{spec:specialization}
Let $Y\to D=\mathbb A^1$ be a smooth morphism of complex algebraic varieties,
and let $g:Y\to B\times D$ be a proper morphism commuting with the maps to
$D$.  Suppose that $Y|_{D^*}$ is
isomorphic to $Y_1\times D^*$ over $D^*$ and that the actual restriction maps
\[
 r_i:H^*(Y,\mathbb Q)\longrightarrow H^*(Y_i,\mathbb Q),
 \qquad i=0,1,
\]
are isomorphisms.  An automorphism $T$ of $Y|_{D^*}$ over $B\times D^*$
then induces an endomorphism $\tau_0$ of $Rg_{0,*}\mathbb Q_{Y_0}$ whose
cohomological action is
\[
 H^*(\tau_0)=r_0r_1^{-1}T_1^*r_1r_0^{-1}.
\]
Neither $B$ nor $Y_i$ is required to be proper.
\end{lemma}

\begin{proof}
Smoothness gives the canonical isomorphism
$\mathbb Q_{Y_0}\simeq\psi\mathbb Q_Y$.
Compatibility of nearby cycles with proper direct image gives
\[
 Rg_{0,*}\mathbb Q_{Y_0}
 \simeq Rg_{0,*}\psi\mathbb Q_Y
 \simeq\psi(Rg_*\mathbb Q_Y).
\]
Apply nearby cycles to the pullback endomorphism induced by $T$ on the
punctured direct image, and use these isomorphisms to define $\tau_0$.

We verify its action on global cohomology without interchanging nearby
cycles with a nonproper direct image to a point.  Put $K=Rg_*\mathbb Q_Y$
on $B\times D$.  Let $\widetilde D^*\to D^*$ be the universal cover and write
\[
 \widetilde j_B:B\times\widetilde D^*\longrightarrow B\times D,
 \qquad i_B:B\times\{0\}\hookrightarrow B\times D,
 \qquad \widetilde K=\widetilde j_B^*K.
\]
Set $\widetilde Y^*=Y|_{D^*}\times_{D^*}\widetilde D^*$.
Proper base change identifies $\widetilde K$ with the direct image of the
constant sheaf from $\widetilde Y^*$.  The nearby-cycle formula downstairs
is $\psi K=i_B^*R\widetilde j_{B,*}\widetilde K$.  Restriction of this
complex gives a natural morphism
\begin{equation}\label{spec:comparison-map}
 \begin{split}
 R\Gamma(\widetilde Y^*,\mathbb Q)
 &\simeq R\Gamma(B\times D,R\widetilde j_{B,*}\widetilde K)\\
 &\longrightarrow R\Gamma(B,\psi K)
 \simeq R\Gamma(Y_0,\mathbb Q).
 \end{split}
\end{equation}
The last isomorphism uses proper nearby-cycle compatibility and smoothness
of $Y\to D$.  The composite of \eqref{spec:comparison-map} with restriction
from $R\Gamma(Y,\mathbb Q)$ is ordinary restriction to $Y_0$.
Indeed the specialization unit for $K$ corresponds, under proper direct
image, to the unit $\mathbb Q_{Y_0}\to\psi\mathbb Q_Y$, which is the
canonical smooth isomorphism used above.

Choose a lift of $1\in D^*$.  The assumed product and the contractibility of
$\widetilde D^*$ identify $H^*(\widetilde Y^*)$ with $H^*(Y_1)$ by actual
restriction to that lift.  Pullback from $H^*(Y)$ therefore becomes $r_1$
and is an isomorphism.  Since its composite with
\eqref{spec:comparison-map} is $r_0$, the latter map is also an isomorphism
on cohomology.

The comparison was constructed on $B\times D$ so that it is equivariant
for the required operator: because $T$ is over $B\times D^*$, it induces
an endomorphism of $K|_{B\times D^*}$ and hence of
$\widetilde K$, $R\widetilde j_{B,*}\widetilde K$, and $\psi K$.
Every arrow in \eqref{spec:comparison-map} respects these actions.
Restricting at the chosen lift identifies the action on its domain with
$T_1^*$, and its action on the codomain is $H^*(\tau_0)$.
Thus the displayed formula for $H^*(\tau_0)$ follows.  No extension of $T$
to an automorphism of $Y$ is asserted or needed.
\end{proof}

Apply Lemma~\ref{spec:specialization} to \eqref{spec:hitchin-line}.  The
resulting endomorphism $\tau_0$ of $Rg_{0,*}\mathbb Q_{Y_{0,0}}$ acts on
cohomology as the translation transported to $Y_{0,0}$ by
Corollary~\ref{log:restriction}.  In particular its action is the operator
whose logarithm is $N_{0,0}$ in Section~\ref{sec:translation}.

A finite polynomial suffices to realize this logarithm over the Hitchin
base.  Choose $M$ as in Lemma~\ref{trans:logarithm}, and choose an integer
$r\ge1$ for which $(\mathsf T^M-1)^r=0$ on total cohomology.  Put
\[
 p(t)=\frac1M\sum_{a=1}^{r-1}
       \frac{(-1)^{a+1}}{a}(t^M-1)^a.
\]
The polynomial $p(\tau_0)$ is a legitimate endomorphism of
$Rg_{0,*}\mathbb Q_{Y_{0,0}}$, and its action on cohomology is $N_{0,0}$.
No unipotence assertion about $\tau_0$ in the derived category is needed.

\subsection{The derived shift}

We have now realized the middle operator in $f'$ over the Hitchin base.
We next include the flag correspondence and the moving-point integration,
keeping every shift explicit.

\begin{proposition}\label{spec:perverse-lowering}
The operator $f'$ of \eqref{trans:fprime} is induced by a morphism
\begin{equation}\label{spec:derived-lowering}
 Rh_*\mathbb Q_X\longrightarrow Rh_*\mathbb Q_X[-2]
\end{equation}
in the constructible derived category on $A$.  Consequently
\[
 f'(P_kH^m(X,\mathbb Q))\subseteq P_{k-2}H^{m-2}(X,\mathbb Q)
\]
for all integers $m,k$.
\end{proposition}

\begin{proof}
Let $O_0=S\times X$ using the scalar product identification of
Lemma~\ref{log:scalar-product}, and write $q=\operatorname{id}_S\times h$.
Both maps in the scalar flag correspondence of
Figure~\ref{log:scalar-correspondence} at $\lambda=0$ are over $\mathcal E$.
The flag projection has relative complex dimension $b$, and the scalar
inclusion has complex codimension $b$.  Thus their pullbacks and proper
Gysin maps give, with
\[
 \begin{gathered}
 F_O=j_*Rq_*\mathbb Q_{O_0},\qquad
 F_Y=Rg_{0,*}\mathbb Q_{Y_{0,0}},\\
 F_Z=R(g_0\circ\iota)_*\mathbb Q_{Z_0},
 \end{gathered}
\]
morphisms
\[
 \mathsf B_0:F_O\longrightarrow F_Y[2b],
 \qquad
 \mathsf A_0:F_Y\longrightarrow F_O[-2b].
\]
Composing these with the appropriate shifts of $p(\tau_0)^2$ gives an
endomorphism of $F_O$ with zero net shift.  Its cohomological action is
$\mathsf A_0N_{0,0}^2\mathsf B_0$.
Since direct image by the closed embedding $j$ is fully faithful, this is
an endomorphism of $Rq_*\mathbb Q_{O_0}$ over $S\times A$.

Let $p_S:S\times A\to A$ be the projection, and write
$K=Rh_*\mathbb Q_X$.  The product identification gives
\[
 Rp_{S,*}Rq_*\mathbb Q_{O_0}
 \simeq K\otimes R\Gamma(S,\mathbb Q).
\]
Push the preceding endomorphism to $A$, precompose with the projection unit
$K\to K\otimes R\Gamma(S,\mathbb Q)$, and postcompose with the orientation
trace
\[
 K\otimes R\Gamma(S,\mathbb Q)\longrightarrow K[-2].
\]
This constructs \eqref{spec:derived-lowering}.  On ordinary cohomology the
unit is $\operatorname{pr}_0^*$, the trace is
$\operatorname{pr}_{0,*}$, and the middle map is the flag composite already
identified.  Hence the resulting operator is exactly $f'$.

For completeness, set $K_{\mathrm{norm}}=K[R]$, the normalization used in
the statement of the theorem.  Functoriality of perverse truncation applied
to \eqref{spec:derived-lowering} gives
\[
 {}^p\tau_{\le k}K_{\mathrm{norm}}
 \longrightarrow{}^p\tau_{\le k}(K_{\mathrm{norm}}[-2])
 =({}^p\tau_{\le k-2}K_{\mathrm{norm}})[-2].
\]
Taking hypercohomology in degree $m-R$ and then its image in ordinary
cohomology proves the stated inclusion.  This uses neither a splitting nor
multiplicativity of the perverse filtration.
\end{proof}

We finish by recording the derived morphisms and their shifts. Retain the
complexes $F_O,F_Z,F_Y$ and $K$ from the proof, and put
$K_S=K\otimes R\Gamma(S,\mathbb Q)$.

The first panel of Table~\ref{spec:shift-table} consists of morphisms on
$\mathcal E$. Compose the flag maps and $p(\tau_0)^2$ there to obtain an
endomorphism of $F_O$, descend it through $j$, and then apply $Rp_{S,*}$.
This produces the degree-zero endomorphism of $K_S$ in the second panel,
whose base is $A$. The unit and trace then give $K\to K[-2]$; it is this
complete morphism that yields the normalized perverse bound above.
No perverse exactness of $Rp_{S,*}$ is used.
The weight columns refer to the cohomological actions after transport
to $\lambda=1$ and use the diagram weights of Section~\ref{sec:diagram}.
In particular $p(\tau_0)$ acts as $N_{1,0}$ after this comparison, while
the unit and trace preserve weight because the cochains of $S$ have
weight zero.

\begin{table}[htbp]
\centering
\small
\begin{tabular}{llrr}
\toprule
\multicolumn{4}{c}{Morphisms on $\mathcal E$}\\
\midrule
Operation & Morphism & Degree & $\Delta W$\\
\midrule
$\pi_0^*$ & $F_O\to F_Z$ & $0$ & $0$\\
$\iota_{0,*}$ & $F_Z\to F_Y[2b]$ & $2b$ & $2b$\\
$p(\tau_0)$ & $F_Y\to F_Y$ & $0$ & $-2$\\
$p(\tau_0)^2$ & $F_Y\to F_Y$ & $0$ & $-4$\\
$\iota_0^*$ & $F_Y\to F_Z$ & $0$ & $0$\\
$\pi_{0,*}$ & $F_Z\to F_O[-2b]$ & $-2b$ & $-2b$\\
Flag composite & $F_O\to F_O$ & $0$ & $-4$\\
\midrule
\multicolumn{4}{c}{Morphisms on $A$}\\
\midrule
Operation & Morphism & Degree & $\Delta W$\\
\midrule
Projection unit & $K\to K_S$ & $0$ & $0$\\
Pushed composite & $K_S\to K_S$ & $0$ & $-4$\\
Orientation trace & $K_S\to K[-2]$ & $-2$ & $0$\\
Complete $f'$ & $K\to K[-2]$ & $-2$ & $-4$\\
\bottomrule
\end{tabular}
\caption{Derived degrees and diagram-weight bounds for the correspondence
maps. A weight bound $s$ means that $W_k$ maps into $W_{k+s}$ after
transport to the Betti side. The row for
$p(\tau_0)$ records $N$ separately; the flag composite uses its square.}
\label{spec:shift-table}
\end{table}
